\section{Descripción del Proyecto}

El Centro Boliviano Americano (CBA) es una institución dedicada a la enseñanza del idioma inglés en la ciudad de Tarija, Bolivia. Actualmente, el proceso de evaluación de nuevos estudiantes se realiza de forma manual, lo que implica tiempo y recursos administrativos significativos.

El presente proyecto consiste en el desarrollo de una aplicación web que permita administrar y rendir exámenes de colocación (\textit{placement tests}) de forma automatizada. El sistema determina el nivel de inglés del estudiante de manera inmediata, basándose en sus respuestas y en una configuración predefinida de niveles y puntajes.

\subsection{Institución Beneficiaria}

\begin{itemize}
    \item \textbf{Nombre:} Centro Boliviano Americano (CBA)
    \item \textbf{Área Responsable:} Área Académica
    \item \textbf{Supervisor Técnico:} Área de Sistemas
    \item \textbf{Supervisor del Proyecto:} Ing. Omar Mancilla
\end{itemize}

\section{Objetivos}

\subsection{Objetivo General}

Desarrollar una aplicación web que permita administrar y rendir exámenes de colocación para determinar automáticamente el nivel de inglés de nuevos estudiantes del Centro Boliviano Americano.

\subsection{Objetivos Específicos}

\begin{itemize}
    \item Diseñar e implementar una base de datos relacional que cumpla con la Tercera Forma Normal (3FN) para almacenar la información de estudiantes, administradores, preguntas, exámenes y resultados.
    \item Implementar un módulo de autenticación que permita el ingreso al sistema mediante CI (Carnet de Identidad) o correo electrónico.
    \item Desarrollar un módulo de administración para la gestión completa de preguntas, niveles y configuración del examen.
    \item Implementar un sistema de generación aleatoria de exámenes con tiempo límite configurable.
    \item Calcular automáticamente el nivel de inglés del estudiante en función del puntaje obtenido.
    \item Generar reportes descargables en formato CSV y Excel.
    \item Asegurar que los resultados históricos no puedan ser modificados.
\end{itemize}

\section{Alcance}

\subsection{Incluye}

\begin{itemize}
    \item Módulo Estudiante: registro, inicio de sesión, rendir examen, ver resultado.
    \item Módulo Administrador: gestión de preguntas, niveles, configuración, dashboard, reportes.
    \item Base de datos relacional con funciones SQL, triggers y Row Level Security.
    \item Autenticación mediante Supabase Auth con soporte para CI o email.
    \item Documentación técnica y manual de usuario.
\end{itemize}

\subsection{No Incluye}

\begin{itemize}
    \item Sistemas de pagos o facturación.
    \item Integraciones con sistemas externos (SSO institucional, LMS).
    \item Aplicación móvil nativa (la web es responsive).
\end{itemize}

\section{Tecnologías Utilizadas}

\begin{table}[H]
\centering
\begin{tabularx}{\textwidth}{l Y}
\toprule
\textbf{Tecnología} & \textbf{Propósito} \\
\midrule
React 19 + TypeScript & Framework de frontend \\
Vite 8 & Build tool y dev server \\
Tailwind CSS 4 & Estilos y diseño responsive \\
Supabase (BaaS) & Backend como Servicio (PostgreSQL, Auth, RLS) \\
PostgreSQL 15 & Base de datos relacional \\
Supabase SDK (\texttt{@supabase/supabase-js}) & Conexión frontend-backend \\
Git + GitHub & Control de versiones \\
Vercel & Despliegue del frontend \\
LuaLaTeX & Documentación técnica \\
\bottomrule
\end{tabularx}
\caption{Tecnologías del proyecto}
\end{table}

\section{Arquitectura General}

El sistema sigue una arquitectura \textbf{Frontend + BaaS} (Backend as a Service). No existe un servidor REST intermedio. El frontend en React consume directamente los servicios de Supabase mediante su SDK oficial.

\begin{figure}[H]
\centering
\begin{lstlisting}[frame=none, numbers=none, basicstyle=\ttfamily\small]
Frontend (React + Supabase SDK)
        |
        v
   Supabase
   ├── PostgreSQL (datos + RLS)
   ├── Auth (autenticación)
   └── Edge Functions (lógica adicional)
\end{lstlisting}
\caption{Diagrama de arquitectura general}
\end{figure}

\section{Reglas de Negocio}

\begin{table}[H]
\centering
\begin{tabularx}{\textwidth}{l Y}
\toprule
\textbf{Regla} & \textbf{Implementación} \\
\midrule
Un estudiante solo puede rendir un examen por día & Trigger SQL \\
El examen tiene tiempo límite configurable & Tabla \texttt{exam\_config} \\
Las preguntas son aleatorias por examen & Función SQL \\
El nivel se calcula automáticamente & Función SQL \\
Los resultados históricos no pueden modificarse & Trigger SQL \\
Validaciones críticas en backend & RLS + Constraints \\
\bottomrule
\end{tabularx}
\caption{Reglas de negocio y su implementación}
\end{table}
